\documentclass[conference]{IEEEtran}
\IEEEoverridecommandlockouts
\usepackage{ctex}
\usepackage{cite}
\usepackage{amsmath,amssymb,amsfonts}
\usepackage{algorithm}
\usepackage{algorithmic}
\usepackage{graphicx}
\usepackage{textcomp}
\usepackage{xcolor}
\usepackage{booktabs}
\usepackage{array}
\usepackage{url}

\def\BibTeX{{\rm B\kern-.05em{\sc i\kern-.025em b}\kern-.08em
		T\kern-.1667em\lower.7ex\hbox{E}\kern-.125emX}}

\begin{document}
	
	\title{面向低空空域的无人机认证前移与分层恶意请求过滤跨域认证及不可洗白撤销方案}
	
	\author{\IEEEauthorblockN{姓名}
		\IEEEauthorblockA{\textit{机构部门}\\
			\textit{机构}\\
			城市, 国家}
	}
	
	\maketitle
	
	\begin{abstract}
		针对低空无人机跨管理域飞行过程中存在的认证触发滞后、跨域信任异构、匿名身份撤销困难以及恶意设备更换软件标识后重新注册等问题，本文提出一种面向低空空域的认证前移与分层恶意请求过滤跨域认证及不可洗白撤销方案。方案采用“低频注册重、在线认证轻”的设计原则，将硬件可信绑定、后量子匿名凭证签发等相对高成本操作集中于注册阶段，并在在线阶段仅使用预生成凭证、短签名和轻量状态查询。针对跨域认证时机，本文设计基于飞行速度、边界法向速度、通信时延和定位误差的弹性边界模型，在无人机实际越界前触发预认证。针对恶意请求处理，构建协议新鲜度、凭证状态、上下文信任评分和布隆过滤器四级过滤链：前两级以极低计算成本丢弃格式异常、重放和已失效凭证，第三级根据历史行为与当前上下文识别高风险请求，最后利用布隆过滤器快速查询已知撤销集合；布隆过滤器命中后再通过稀疏默克尔树进行精确状态确认，避免将概率型数据结构误作为最终撤销依据。正式认证采用挑战--响应和状态证明完成双向认证，并建立短期会话密钥。设备确认恶意后，由联盟链智能合约记录撤销事件，同时更新SMT状态和硬件PUF死名单；后续重新注册必须通过硬件死名单检查，从物理设备层面阻断身份洗白。本文进一步采用IPFS保存大规模状态数据、联盟链锚定状态根，并以BLS短签名降低高密度场景通信开销。方案形成“注册--认证前移--分层过滤--正式认证--行为监测--精确撤销--硬件防洗白”的完整安全闭环，为长距离巡检、物流配送和跨域安防等行业级无人机提供一种可实现的安全认证框架。
	\end{abstract}
	
	\begin{IEEEkeywords}
		无人机跨域认证；认证前移；恶意请求过滤；布隆过滤器；稀疏默克尔树；PUF；区块链；隐私保护
	\end{IEEEkeywords}
	
	\section{引言}
	
	\subsection{研究背景}
	近年来，低空空域逐渐由传统的单机、单域应用向多主体协同和跨区域连续运行转变。无人机在电力巡检、交通巡查、应急救援、物流运输以及低空安防等任务中通常需要跨越多个管理域，而不同管理域在身份注册、访问控制和安全策略方面具有独立性。无人机一旦进入新的管理域，必须在较短时间内完成身份确认和权限迁移，否则容易出现通信中断、任务暂停甚至飞行安全风险。因此，跨域认证的核心问题已经由“能否认证”进一步转变为“能否在正确的时间、以可接受的代价完成认证”。
	
	从网络安全角度看，低空无人机系统还具有明显的开放无线环境特征。无线链路容易遭受窃听、重放、伪造和洪泛攻击，攻击者可以持续发送低成本请求，使边缘认证节点将大量计算资源消耗在无效请求上。传统认证协议通常首先执行签名验证、证书验证或复杂的密钥协商，再判断请求是否具有进一步处理的必要性，这种处理顺序在高并发场景下容易形成“高成本入口”。如果攻击者掌握合法设备凭证，其请求甚至能够通过最基本的身份检查，传统基于凭证有效性的认证机制难以进一步区分正常设备与恶意设备。
	
	此外，匿名认证与设备撤销之间存在天然矛盾。为了保护无人机身份隐私，系统通常不会在无线消息中直接传输真实身份，而是使用假名、临时标识或匿名凭证。然而，一旦匿名标识被撤销，系统仍然需要判断该设备后续生成的新标识是否属于同一物理设备。若撤销机制仅绑定软件身份，则攻击者可以通过修改设备标识、重新生成密钥或更换应用层凭证逃避撤销。因此，低空无人机安全体系需要同时解决“通信身份匿名”和“物理身份连续”两个看似相反的问题。
	
	基于上述特点，本文将跨域认证过程划分为低频注册阶段和高频在线阶段，并进一步将恶意请求处理前移到正式认证之前。该设计并非简单增加密码算法数量，而是通过合理的功能分层改变系统的计算路径：在设备注册时承担较高成本，在实际飞行期间尽可能减少重复计算；对于进入系统的请求，则按照由轻到重的顺序逐层筛选，使明显异常请求在到达高成本密码认证模块之前被丢弃。
	随着低空经济和无人机规模化应用的发展，行业级无人机逐渐承担跨行政区域巡检、物流运输、应急通信和全域安防等任务。与单一管理域内飞行相比，跨域飞行要求无人机在不同管理机构之间连续完成身份确认、权限验证和状态检查。传统集中式PKI体系依赖单一认证中心，存在单点故障和跨域信任链较长的问题；当大量无人机同时接近域边界时，集中认证还可能产生明显的计算和通信拥塞。
	
	另一方面，匿名认证虽然能够降低飞行轨迹被长期关联的风险，却使恶意设备撤销更加困难。攻击者在获得合法凭证后可能实施拒绝管控、异常航线飞行或伪造控制消息。当软件ID、MAC地址或固件标识均可修改时，仅撤销软件层假名并不能阻止同一物理设备重新注册。与此同时，传统CRL或频繁链上查询难以同时满足低空场景中的实时性和大规模设备状态管理需求。
	
	因此，低空无人机跨域认证并非单纯的“提高密码强度”问题，而是一个由认证时机、请求筛选、状态证明、动态撤销和硬件身份连续性共同组成的系统问题。本文围绕这一问题建立统一协议，使各密码和数据结构组件只承担明确的单一职责。
	
	\subsection{研究问题}
	围绕低空跨域飞行场景，本文重点研究认证时机、认证成本、恶意请求处理和设备撤销之间的协同关系。首先，认证触发点需要与无人机运动状态和网络状态相适配。固定距离边界虽然容易实现，但无法反映不同飞行速度、通信时延和定位误差对认证准备时间的影响。其次，注册和在线认证具有明显不同的频率特征，若所有安全操作均放置在在线阶段，会导致无人机在高并发和快速移动情况下承担不必要的计算负担。
	
	另一方面，恶意设备可能同时表现为两种形式：一类是没有合法凭证的外部伪造请求，另一类是具有合法凭证但在任务过程中产生异常行为的内部恶意设备。前者适合采用格式、新鲜度和状态检查快速过滤，后者则需要结合历史行为和飞行上下文进行判断。因此，单一过滤器难以覆盖完整攻击面，需要建立层次化的请求处理路径。
	
	最后，撤销结果必须同时满足实时查询和状态可信两个条件。Bloom Filter可以以较低空间成本完成高速集合判断，但它本质上是概率型结构；区块链能够提供不可篡改的撤销记录，但直接进行链上查询会带来额外时延；SMT能够提供高效状态证明，却不适合直接承担大规模请求入口过滤。本文因此将三者分别置于快速过滤、精确证明和可信锚定三个不同位置，形成明确的职责边界。
	本文主要解决以下问题：
	\begin{enumerate}
		\item 如何在无人机真正越界之前完成跨域准备，避免认证时延导致越界后短暂失管；
		\item 如何将高成本操作集中于低频注册阶段，同时降低在线跨域认证负担；
		\item 如何在正式认证之前利用多级轻量过滤机制尽早丢弃明显恶意请求；
		\item 如何实现“快速撤销集合查询”和“精确撤销状态证明”的职责分离；
		\item 如何防止被撤销设备通过修改软件身份重新注册，实现硬件级不可洗白；
		\item 如何在匿名性、可撤销性和跨域审计之间建立平衡。
	\end{enumerate}
	
	\subsection{主要贡献}
	本文的贡献并不局限于引入新的单一密码原语，而是针对低空无人机跨域认证的时序特点，对认证、过滤和撤销流程进行重新组织。具体而言，本文首先建立一种与飞行运动状态相关的认证前移模型，使认证准备时间成为边界触发决策的一部分；其次，将较重的设备可信绑定和匿名凭证建立操作集中于低频注册阶段，并将在线阶段压缩为凭证使用和轻量状态检查；再次，提出四级过滤链，在不同安全判断阶段使用不同复杂度的检查方式；最后，通过PUF硬件根、SMT状态和区块链撤销记录形成从软件身份到物理设备的连续撤销关系。
	本文贡献概括如下：
	\begin{enumerate}
		\item 提出一种弹性边界驱动的认证前移机制，根据无人机接近边界所需时间和系统认证准备时间动态确定预认证触发区间。
		\item 提出“重型注册、轻量在线”的混合密码架构，在注册阶段完成PUF绑定和后量子匿名凭证建立，在线阶段避免重复执行高成本注册操作。
		\item 设计四级恶意请求过滤链，依次执行协议/新鲜度过滤、凭证状态过滤、上下文信任过滤和Bloom Filter集合过滤；Bloom Filter只负责快速候选筛选，最终撤销结论由SMT状态证明给出。
		\item 构建基于PUF死名单和链上撤销记录的不可洗白机制，使软件层假名更新无法绕过物理设备级撤销。
		\item 建立联盟链与IPFS协同状态管理架构，将链上不可篡改锚定与链下大规模状态存储分离，降低状态同步压力。
	\end{enumerate}
	
	\section{相关工作}
	
	\subsection{无人机跨域认证}
	无人机跨域认证的主要目标是在不共享全部内部身份数据库的情况下，使目标管理域能够确认访问设备及其所属域的合法性。现有方案通常利用联盟链、边缘计算、群签名、匿名凭证或轻量密钥协商降低跨域认证成本。对于低空场景而言，跨域认证还必须考虑无人机高速移动带来的时间约束：如果认证仅在实际越界后开始，即使密码协议本身只需要较短时间，网络排队、域间通信和状态同步也可能使认证结果晚于业务切换时刻。
	
	因此，本文不将跨域认证视为一个孤立的握手过程，而是把它划分为“准备”和“授权”两个阶段。预认证阶段主要完成目标域授权上下文建立和必要状态准备，正式认证阶段才决定是否允许设备获得目标域服务权限。这种分离能够降低实际越界瞬间的认证压力，并为后续恶意请求过滤提供一个独立入口。
	现有无人机跨域认证研究主要围绕区块链、信任管理、匿名认证和轻量密码等方向展开。Rui等提出基于信任评价的UAV跨域批量认证机制，提高批量接入效率，但未重点解决后量子安全和硬件级撤销问题\cite{b1}。Qiao等将区块链与轨迹隐私结合，用于无人机跨域认证和飞行轨迹保护\cite{b2}。Feng等进一步研究可信高效的区块链UAV跨域认证协议\cite{b21}。这些工作说明区块链可以降低跨域信任依赖，但高频链上操作仍可能影响实时性。
	
	\subsection{可撤销匿名认证}
	匿名认证通常通过假名或匿名凭证隐藏真实身份，使观察者难以从网络报文直接获得无人机的长期身份。然而，匿名性越强，撤销过程所面临的状态管理问题越突出。传统证书撤销列表需要周期性广播，随着设备数量增加会产生较大的通信和存储开销；链上逐设备记录虽然能够增强不可篡改性，但不适合将所有在线查询直接转换为链上交易。
	
	本文采用分层状态管理方式。在线通信只使用匿名PID，目标域通过本地缓存和Bloom Filter完成快速判断；需要进行安全确认时，再利用SMT证明设备当前状态；撤销事件最终由联盟链记录。由此，匿名身份的使用范围、撤销状态的证明方式以及不可篡改记录分别由不同机制承担，避免单一数据结构承担过多职责。
	车联网和工业物联网中的可撤销匿名认证为无人机系统提供了重要参考。Li等研究区块链辅助的可撤销跨域认证\cite{b3}，Zeng等研究可撤销跨域匿名认证\cite{b17}。然而，仅依赖软件假名或证书状态的撤销机制不能从根本上解决设备更换身份后的重新注册问题。本文因此将匿名假名和物理PUF身份分离：正常通信使用匿名PID，而撤销和重新注册审计使用硬件绑定状态。
	
	\subsection{恶意请求过滤与信任评估}
	恶意请求过滤的关键并非简单提高检测模型复杂度，而是控制不同类型请求在系统中的“计算路径”。对于明显格式错误、时间戳过期或Nonce重复的请求，继续执行密码验证没有安全收益；对于已经确认失效的凭证，同样没有必要进入行为评分；只有无法通过低成本规则直接判定的请求，才值得投入更多计算资源。因此，过滤机制应遵循“低成本检查优先，高成本判断后置”的原则。
	
	Bloom Filter适合这一过程中的集合筛选，但不适合独立完成安全决策。其假阳性意味着某个合法设备可能被暂时标记为“可能撤销”，但只要命中后执行精确状态验证即可消除这一问题。相比之下，信任评分用于处理合法凭证持有者的行为异常，它回答的是“该请求是否值得进一步审查”，而不是“该设备是否已经被正式撤销”。这种语义分离是本文过滤链设计的基础。
	恶意请求过滤通常采用规则检测、速率限制、信誉评价、异常检测和概率型集合数据结构。Bloom Filter具有低空间开销和常数级近似集合查询优势，适合部署在高频请求入口；但其存在假阳性，因此不适合直接承担最终安全判决。信任管理研究则能够利用历史行为、当前行为和上下文信息识别具有合法凭证但行为异常的节点。本文将两类机制组合起来，使信任评分负责“可疑性判断”，Bloom Filter负责“已知撤销集合快速查询”，SMT负责最终精确状态确认。
	
	\subsection{研究空白}
	现有研究在跨域认证、匿名保护、区块链撤销和无人机信任管理方面已经形成较为丰富的成果，但这些机制往往按照单一安全目标分别设计，较少从完整飞行生命周期出发讨论认证触发、恶意请求处理和物理设备撤销之间的连续关系。尤其是在合法凭证被恶意使用的情况下，仅依赖密码学认证无法回答设备是否正在实施异常行为；而仅依赖行为检测又缺乏可信的撤销状态传播机制。
	
	本文的设计重点由此转向系统级协同：利用弹性边界决定“何时准备认证”，利用预认证机制决定“如何提前准备”，利用四级过滤决定“哪些请求值得进入正式认证”，利用SMT决定“撤销状态是否真实”，再利用PUF死名单回答“被撤销的物理设备是否可以重新获得身份”。这一链路构成本文与单纯跨域认证协议的主要区别。
	综合已有研究可发现，当前方案通常分别优化认证、隐私或撤销中的某一个环节，而缺少从边界触发到恶意设备最终撤销的完整闭环。本文的核心思路不是堆叠密码原语，而是按照系统生命周期重新分配功能，使每一模块仅执行一种核心任务。
	
	\section{系统模型与安全需求}
	
	\subsection{系统实体}
	无人机$U_j$是系统中的移动终端，负责采集任务数据并通过无线链路访问不同管理域。母域$\mathcal{D}_A$负责设备注册、初始授权以及设备真实身份的受控映射。目标域$\mathcal{D}_B$负责接收预认证请求、执行过滤和正式认证。域边缘节点$RSU_B$部署过滤缓存、信任评分和Bloom Filter，以缩短请求处理路径。联盟链网络$\mathcal{BC}$由各管理域共同维护，用于保存撤销事件、状态根和版本信息；IPFS用于存储较大规模的状态树和审计证据。
	
	这种实体划分使系统中的计算位置具有明确边界。无人机主要执行注册所得凭证的使用和轻量密码操作；边缘节点承担高频过滤；管理域负责策略判断和授权；联盟链承担最终可信记录。因而即使请求数量突然增加，也不会迫使每一条无线请求都直接访问链上状态。
	系统包含无人机$U_j$、母域$\mathcal{D}_A$、目标域$\mathcal{D}_B$、域边缘节点$RSU_B$、联盟链网络$\mathcal{BC}$和IPFS网络。母域负责设备注册和初始授权，目标域负责跨域接入验证，联盟链负责跨域可信状态锚定，IPFS保存完整SMT和硬件死名单等大规模数据。
	
	\subsection{信任关系}
	不同管理域之间不要求共享全部设备数据库，而是通过联盟链建立最小必要的信任锚。域公钥、撤销状态根和版本信息作为跨域验证所需的公共状态，由联盟链提供一致性保证。母域掌握设备真实身份与匿名PID之间的受控映射，目标域只获得完成当前任务所需的最小信息。
	
	这种最小披露原则能够降低跨域共享造成的隐私扩散。对于正常业务，目标域只需要确认匿名凭证、当前状态和授权范围；只有发生安全事件并满足审计条件时，才由授权实体执行匿名身份的反向解析。
	各管理域作为联盟链成员节点运行共识协议。域间公钥和状态根通过智能合约进行认证。无人机不需要保存其他所有域的完整设备列表，只保存自身注册凭证、匿名PID、PUF派生秘密以及必要的缓存状态。
	
	\subsection{攻击者模型}
	本文考虑具有较强网络控制能力但无法直接破坏联盟链共识和PUF硬件根的攻击者。攻击者能够监听公开无线通信、复制和重放历史消息、修改消息字段、构造伪造请求，也可能控制已经获得合法凭证的无人机并主动产生异常行为。对于被撤销设备，攻击者进一步尝试修改RealID、MAC地址、软件密钥或重新申请匿名PID。
	
	此外，本文考虑请求洪泛攻击。攻击者不一定需要伪造出有效身份，只需持续发送大量低成本请求即可消耗边缘节点的签名验证和状态查询资源。因此过滤链需要在尽可能靠近入口的位置运行，并在高成本密码运算之前完成大量低成本拒绝。对于Bloom Filter的假阳性，攻击者不能利用该性质获得错误放行，因为命中结果必须经过SMT精确核验。
	攻击者可窃听、重放、篡改和伪造无线请求；可控制已经获得合法软件凭证的恶意无人机；可尝试修改设备ID、MAC、固件和随机注册参数；可利用撤销状态同步延迟重复接入。攻击者不能直接提取PUF内部不可导出的硬件响应，也不能破坏联盟链的基本共识安全假设。
	
	\subsection{安全需求}
	除传统的机密性、完整性和认证性外，低空无人机系统还需要满足时序安全和状态连续性要求。时序安全要求认证能够在预期跨域时间窗口内完成；状态连续性要求设备撤销状态能够从母域传播到目标域，并在重新注册时继续生效。匿名性要求无线观察者不能直接获得真实设备身份，不可链接性则要求不同任务会话之间尽量避免长期标识关联。
	
	对于过滤链，安全需求还包括单调安全性：增加过滤层不应降低最终认证的安全级别。即任何过滤器的“通过”都不能替代正式认证，任何概率型过滤器的“命中”也不能直接等价于恶意判决。最终授权必须由密码学凭证和可信状态共同决定。
	方案需要满足双向认证、消息不可伪造、重放抵抗、匿名性、不可链接性、动态撤销、跨域状态一致性和硬件级防洗白等要求。对于Bloom Filter造成的假阳性，系统应保证其只导致额外核验而不导致错误放行；因此最终安全判决不能由Bloom Filter单独给出。
	
	\section{系统总体架构}
	系统总体架构按照设备生命周期组织，而不是按照密码算法组织。注册阶段建立长期信任根；认证前移阶段根据飞行轨迹预测触发时刻；预认证阶段建立目标域上下文；过滤阶段将请求按照风险和撤销状态逐层筛选；正式认证阶段完成最终授权；任务执行阶段持续更新信任状态；确认恶意后执行正式撤销并同步硬件死名单。
	
	在该架构下，各模块之间形成严格的前后依赖关系。任何设备只有先获得有效注册凭证，才能参与预认证；预认证请求必须经过过滤，才能进入正式认证；正式撤销必须在精确状态确认后执行；撤销后的设备即使生成新的软件身份，也必须再次经过PUF死名单检查。因而系统不存在“跳过中间环节直接恢复身份”的旁路。
	
	从部署位置看，系统采用终端、边缘、管理域和联盟链四级协同结构。无人机承担与任务直接相关的轻量运算，边缘节点负责高频过滤和缓存访问，管理域负责安全策略和异常审计，联盟链提供跨域状态一致性。该结构能够避免将高频无线请求全部转发到中心节点，也避免把完整撤销数据库直接复制到每一架无人机。
	
	系统遵循如下生命周期：
	\begin{equation}
		Registration\rightarrow PreMove\rightarrow Filtering\rightarrow
		Authentication\rightarrow Monitoring\rightarrow Revocation\rightarrow AntiWashing.
	\end{equation}
	
	其中，Registration为低频重操作；PreMove负责提前准备；Filtering负责低成本排除明显异常；Authentication完成最终身份和状态确认；Monitoring持续更新行为信任；Revocation负责不可逆状态变更；AntiWashing负责阻止已撤销物理设备重新注册。
	
	\subsection{功能解耦原则}
	协议设计中最容易产生安全分析困难的问题之一，是多个机制共同承担同一个安全目标。例如，如果Bloom Filter、SMT和智能合约均被描述为“撤销机制”，就难以准确说明每一个组件在攻击场景下的必要性。本文因此按照“一个组件一个主要职责”的原则进行设计。
	
	这种解耦不仅有利于理论分析，也便于后续实现和性能测试。实验可以分别测量过滤层耗时、SMT证明耗时、链上确认耗时和PUF注册耗时，而不会因为模块重叠导致指标重复计算。
	为避免多个模块同时完成同一功能，本文作如下限定：
	\begin{itemize}
		\item PUF只负责物理设备绑定，不参与在线会话认证计算；
		\item 匿名PID只负责正常通信中的身份隐私，不直接作为物理身份；
		\item PRE只负责跨域授权密文转换，不承担撤销判定；
		\item Bloom Filter只负责快速集合查询，不承担最终撤销判决；
		\item SMT只负责精确状态证明，不负责行为异常检测；
		\item 区块链只负责可信状态锚定和撤销事件记录，不保存完整设备数据；
		\item IPFS只负责大规模状态数据保存，不承担信任决策。
	\end{itemize}
	
	\section{所提跨域认证方案}
	
	\subsection{系统初始化}
	系统设置安全参数$\lambda$，哈希函数$H:\{0,1\}^{*}\rightarrow\{0,1\}^{256}$，以及SMT根状态、Bloom Filter和联盟链智能合约。各管理域生成后量子签名密钥对：
	\begin{equation}
		(pk_i,sk_i)\leftarrow \mathrm{PQ.KeyGen}(1^\lambda).
	\end{equation}
	域公钥、状态根和版本号写入联盟链。
	
	\subsection{无人机硬件绑定与重型注册}
	与在线认证相比，设备注册通常只在设备首次加入系统、密钥更新或管理域重新授权时发生，其调用频率明显较低。因此本文允许注册阶段承担相对较高的计算和通信成本。该策略尤其适合近年来计算能力和密码加速能力不断提高的行业级无人机：与其在每次跨域切换时重复执行复杂操作，不如在注册阶段预先建立可长期使用的安全材料。
	
	PUF用于建立设备与硬件之间的关联。由于PUF响应受制造差异影响，每次物理测量可能存在少量噪声，因此首先通过模糊提取器得到稳定密钥，而不是直接把原始响应作为密码密钥。真实身份只参与注册承诺计算，不在在线无线认证消息中直接传输。
	
	匿名凭证的签发放置于注册阶段，使无人机后续跨域时能够直接使用已经获得的凭证材料。对于后量子安全，域级长期签名采用标准化后量子签名算法；匿名凭证部分则独立采用具有安全证明的格基盲签名构造。这样既避免把标准化签名算法误认为盲签名，又使系统长期信任与匿名性之间的功能边界保持清晰。
	
	注册阶段是低频操作，因此允许使用较高计算成本换取长期凭证安全。无人机首先从SRAM-PUF读取原始响应$R_j$，通过模糊提取器获得稳定密钥：
	\begin{equation}
		(K_j,Helper_j)\leftarrow FE.Gen(R_j).
	\end{equation}
	
	无人机生成随机秘密$s_j$并计算设备绑定承诺：
	\begin{equation}
		C_j=H(K_j\parallel s_j\parallel RealID_j\parallel t_{reg}).
	\end{equation}
	
	为避免将传统RSA盲签名公式与Dilithium混用，本文将匿名凭证抽象为经过安全证明的格基盲签名接口：
	\begin{equation}
		(\sigma_j,\omega_j)\leftarrow BlindSign_{PQ}(C_j),
	\end{equation}
	其中$\omega_j$表示必要的盲化辅助信息。注册服务器只看到盲化后的注册消息，不能直接获得无人机后续使用的匿名PID。设备解盲后计算：
	\begin{equation}
		PID_j=H(C_j\parallel \sigma_j).
	\end{equation}
	
	母域随后在SMT中建立状态叶：
	\begin{equation}
		Leaf_j=H(PID_j\parallel Status_j\parallel Exp_j\parallel Ver_j).
	\end{equation}
	完整SMT存入IPFS，根哈希$Root_A$及CID写入联盟链。
	
	\subsection{弹性边界与认证前移}
	固定地理距离并不能准确反映认证所需时间。例如，同一架无人机在低速巡航和高速穿越边界时，即使距离边界相同，其可用于认证准备的时间也完全不同；无线链路拥塞和定位误差同样会改变实际可用时间。因此本文将边界触发条件定义为时间裕度问题，并通过法向速度估计无人机到达边界的剩余时间。
	
	当剩余飞行时间接近认证准备时间时，系统提前启动预认证。此时目标域并未立即授予正式飞行权限，而只是准备验证上下文和授权状态。这样既避免认证过早造成无意义的资源占用，也避免认证过晚导致设备已经进入目标域却仍处于未授权状态。
	
	定义无人机当前位置到目标行政边界的最短距离为$L(t)$，速度在边界法向上的分量为$v_n(t)$. 根据系统剩余认证时间估计值$T_{auth}$、网络抖动裕度$T_{jit}$和定位误差$\sigma_{GNSS}$，设置动态安全裕度：
	\begin{equation}
		d(t)=\max(0,v_n(t)(T_{auth}+T_{jit})+3\sigma_{GNSS}).
	\end{equation}
	
	当
	\begin{equation}
		L(t)\le d(t),
	\end{equation}
	无人机进入预认证触发区。若预计跨越边界时间小于预认证所需时间，则立即发送预认证请求。相比固定距离阈值，该方法可以根据速度和通信状态自动扩大或缩小触发范围。
	
	\subsection{跨域预认证}
	预认证的核心作用是把“域间准备”从真正的越界时刻向前移动。请求中包含目标域标识、一次性Nonce和时间戳，使目标域能够确认该请求属于当前飞行过程，而不是历史请求的重复发送。母域通过域间预建立的重加密关系向目标域转换必要授权信息，目标域只接收与本次访问相关的密文和状态。
	
	预认证完成后，目标域仅建立候选接入上下文，不直接授予完整任务权限。正式权限仍取决于后续过滤和正式认证。该设计可以避免攻击者利用预认证接口直接获得服务，同时将复杂的跨域准备工作与最终授权决策分离。
	
	无人机生成一次性随机数$N_j$和时间戳$T_j$：
	\begin{equation}
		M_{pre}=PID_j\parallel D_B\parallel N_j\parallel T_j.
	\end{equation}
	
	利用在线轻量签名生成：
	\begin{equation}
		\sigma_{pre}=BLS.Sign(sk_j,M_{pre}).
	\end{equation}
	
	母域验证请求的新鲜度后，通过预先建立的域间重加密密钥执行：
	\begin{equation}
		CT_B=PRE.ReEnc(rk_{A\rightarrow B},CT_A).
	\end{equation}
	
	目标域获得预认证信息后，仅将$PID_j$加入预备接入缓存，不立即授予正式飞行权限。这样可以把耗时较长的跨域授权准备提前到实际越界之前。
	
	\subsection{四级恶意请求过滤}
	四级过滤按照计算成本递增的顺序组织。第一层只检查协议字段和新鲜度；第二层检查本地凭证状态；第三层根据行为和上下文计算风险；第四层查询已知撤销集合。只有无法在低成本阶段被排除、且未命中撤销集合的请求才进入正式认证。
	
	这一结构的主要优势在于降低高成本接口暴露面。设到达请求总量为$N$，经过四层过滤后进入正式认证的比例为$\rho$，则正式认证模块处理的请求量约为$\rho N$。在攻击流量占比较高的情况下，只要前三层能够有效过滤明显异常请求，就可以显著降低$\rho$。因此过滤链的价值不仅在于识别恶意设备，也在于控制认证计算资源的消耗。
	
	过滤层之间还存在明确的状态传递关系。第一层输出“语法和新鲜度有效”的请求，第二层在此基础上确认凭证仍具有业务有效性，第三层进一步判断行为风险，第四层则针对撤销集合进行概率型快速查询。任何一层都不能直接替代正式认证。
	四级过滤按照计算成本递增的顺序组织。第一层只检查协议字段和新鲜度；第二层检查本地凭证状态；第三层根据行为和上下文计算风险；第四层查询已知撤销集合。只有无法在低成本阶段被排除、且未命中撤销集合的请求才进入正式认证。
	
	这一结构的主要优势在于降低高成本接口暴露面。设到达请求总量为$N$，经过四层过滤后进入正式认证的比例为$\rho$，则正式认证模块处理的请求量约为$\rho N$。在攻击流量占比较高的情况下，只要前三层能够有效过滤明显异常请求，就可以显著降低$\rho$。因此过滤链的价值不仅在于识别恶意设备，也在于控制认证计算资源的消耗。
	
	过滤链位于预认证请求接收之后、正式跨域信任建立之前。其目标不是证明设备绝对恶意，而是以尽可能低的成本排除明显异常请求，并将高风险请求送入精确核验。
	
	\subsubsection{第一层：协议完整性与新鲜度过滤}
	第一层过滤不执行签名验证，而是检查消息是否满足最基本的协议语法和时序约束。该设计具有明显的工程意义：攻击者可以非常低成本地构造缺字段、错误长度或过期时间戳的报文，如果系统仍为每个报文执行密码运算，将形成不必要的资源浪费。
	
	时间窗口与Nonce缓存共同覆盖两类典型重放方式。时间窗口限制消息生命周期，目标域随机挑战保证不同正式认证会话之间的消息不可复用，Nonce缓存则阻止有效窗口内的重复发送。该层的判定结果只决定请求是否具有继续处理的必要性，并不作为最终身份认证依据。
	检查消息长度、字段格式、域标识、时间戳和Nonce：
	\begin{equation}
		|T_{now}-T_j|\le\Delta T.
	\end{equation}
	若Nonce已经出现于短期缓存，直接拒绝。该层不执行昂贵密码运算，主要抵御重放、畸形包和无效请求洪泛。
	
	\subsubsection{第二层：凭证状态过滤}
	通过第一层的请求进入本地凭证状态检查。边缘节点缓存当前有效凭证版本、有效期和状态位，能够在不访问区块链的情况下处理大多数过期和已知失效请求。缓存数据并非系统最终信任源，其版本号和状态根定期与联盟链锚定信息同步。
	
	这种设计将高频读取与低频可信状态更新分离。对于普通合法请求，边缘节点可以直接读取缓存；对于状态存在不一致或版本异常的请求，则升级到精确状态核验。因而链上状态不会成为所有认证请求的必经瓶颈。
	根据本地域缓存查询$PID_j$对应的凭证版本、有效期和撤销状态。若
	\begin{equation}
		Status_j\neq Active
	\end{equation}
	或
	\begin{equation}
		T_{now}>Exp_j,
	\end{equation}
	则直接拒绝。该层避免大量明显失效请求进入后续信任计算。
	
	\subsubsection{第三层：上下文信任过滤}
	第三层针对“身份合法但行为可疑”的设备。系统不依赖单一异常指标，而是综合近期行为、历史信誉和当前飞行上下文形成Trust值。行为维度描述设备最近的请求模式和任务执行状态，历史维度反映长期可靠程度，上下文维度则考虑当前位置、速度、通信质量和任务时段等因素。
	
	信任值采用连续区间而非简单二元判决。低于$\tau_L$的请求可直接拒绝，高于$\tau_H$的请求进入正常流程，中间区间进入Review状态。这样可以减少单次异常导致的误撤销，并为后续安全审计保留足够信息。需要强调的是，信任评分本身不是撤销操作，它只负责决定请求是否需要进一步检查。
	对通过前两层的请求计算轻量风险分数：
	\begin{equation}
		Trust_j=w_bB_j+w_hH_j+w_cC_j,
	\end{equation}
	其中$B_j$为近期行为评分，$H_j$为历史信誉，$C_j$为当前飞行上下文评分，且
	\begin{equation}
		w_b+w_h+w_c=1.
	\end{equation}
	
	行为评分可根据认证失败率、异常请求频率和航线偏离程度计算；历史评分采用指数加权更新：
	\begin{equation}
		H_j^{new}=\alpha H_j^{old}+(1-\alpha)H_j^{cur}.
	\end{equation}
	
	设置低风险阈值$\tau_L$和高风险阈值$\tau_H$：
	\begin{equation}
		Decision_j=
		\begin{cases}
			Reject,&Trust_j<\tau_L,\\
			Review,&\tau_L\le Trust_j<\tau_H,\\
			Pass,&Trust_j\ge\tau_H.
		\end{cases}
	\end{equation}
	为避免误拒绝，处于Review区间的请求进入后续精确状态检查，而非直接判定为恶意。
	
	\subsubsection{第四层：Bloom Filter}
	第四层负责回答一个非常具体的问题：当前匿名PID是否可能出现在已知撤销集合中。与直接遍历撤销列表相比，Bloom Filter能够显著降低集合查询的存储和计算成本，尤其适用于边缘节点需要处理大量请求的场景。
	
	由于Bloom Filter允许假阳性，本文明确规定其返回值只具有“候选状态”语义。当查询结果为0时，可认为该PID不属于当前过滤器维护的撤销集合；当结果为1时，系统不能直接把设备判定为恶意，而应读取对应SMT证明进行精确确认。该设计保证了概率型数据结构不会改变系统的最终安全语义。
	系统维护撤销PID集合对应的Bloom Filter $BF_R$. 对请求执行：
	\begin{equation}
		BF_R.Query(PID_j).
	\end{equation}
	若结果为0，则PID不属于当前已知撤销集合，可以进入正式认证；若结果为1，仅表示“可能属于撤销集合”，必须读取SMT状态进行精确验证。因此Bloom Filter的假阳性只增加验证成本，不降低安全性。
	
	\subsection{正式跨域认证}
	经过四级过滤后，请求进入正式跨域认证。此时进入认证模块的请求数量已经显著低于无线入口请求数量，因此系统可以将计算资源集中于真正具有安全价值的验证过程。目标域生成随机挑战，要求无人机利用注册阶段获得的匿名凭证完成响应；同时，目标域验证母域提供的状态证明，确认该设备在当前版本的全局状态中仍处于Active。
	
	正式认证与预认证之间存在明确的授权边界。预认证表示目标域已经准备好接收设备，正式认证才表示设备获得当前任务的服务权限。即使攻击者绕过过滤器构造合法格式请求，也必须继续满足密码学认证和最新状态证明要求。
	
	目标域向无人机发送随机挑战$R_B$和时间戳$T_B$。无人机生成：
	\begin{equation}
		M_{auth}=PID_j\parallel R_B\parallel T_B\parallel D_B.
	\end{equation}
	
	无人机利用注册阶段获得的匿名凭证对应证明材料生成认证响应。目标域首先验证签名和挑战新鲜度，然后验证SMT状态证明：
	\begin{equation}
		SMT.Verify(Root_A,PID_j,Proof_j)=1.
	\end{equation}
	
	状态叶中包含$Active$标记、有效期和版本号，因此被撤销设备即使仍持有历史匿名凭证，也不能通过最新状态证明。
	
	双方完成验证后采用临时密钥协商生成：
	\begin{equation}
		K_{sess}=KDF(g^{ab}\parallel PID_j\parallel R_B).
	\end{equation}
	$K_{sess}$仅用于当前跨域会话，任务结束后立即失效。
	
	\subsection{持续行为监测与信任更新}
	跨域认证成功后，设备并不会自动获得永久可信状态。任务执行过程中的行为可能发生变化，因此目标域继续以滑动窗口方式记录必要的安全事件。为了减少隐私泄露，系统只保存与安全决策直接相关的统计特征，不要求长期保存完整飞行轨迹作为普通认证数据。
	
	信任值更新采用指数加权方式，使近期行为能够快速影响评分，同时保留一定的历史稳定性。当设备短时间内发生一次偶发异常时，评分不会立即下降到撤销阈值；当异常持续出现时，历史和当前行为项会共同推动评分下降，最终进入审计和撤销流程。
	
	认证成功并不意味着设备永久可信。目标域在任务执行期间持续采集最小必要的安全事件摘要，包括异常请求次数、认证失败次数、航线偏离比例和非法控制指令次数。系统按照滑动窗口更新$Trust_j$。当信任值下降到$\tau_R$以下时触发撤销候选事件，而不是立即执行链上撤销。
	
	\subsection{恶意请求确认与正式撤销}
	本文将“可疑”和“已撤销”严格区分。过滤层可以拒绝请求，但只有在持续异常证据达到策略阈值并经过授权节点确认后，设备才进入正式撤销状态。这种两阶段判决方式可以减少由单次网络异常、GNSS波动或无线丢包造成的误撤销。
	
	正式撤销由智能合约执行并形成不可篡改事件。撤销事件包含匿名PID、原因类别、状态版本和时间信息。SMT随后更新对应叶节点，目标域通过新的状态根确认撤销结果。由于链上记录、状态树和边缘缓存具有版本关系，目标域可以检测到旧状态，从而避免因缓存尚未更新而长期接受已撤销设备。
	
	当请求通过过滤链但仍表现出持续异常行为时，进入精确审计阶段。审计节点重新验证：
	\begin{equation}
		Audit=(Evidence,Trust_j,StateProof,DomainSignature).
	\end{equation}
	
	只有满足策略阈值并获得授权管理节点确认后，智能合约执行撤销：
	\begin{equation}
		Revoke(PID_j,Reason,Version,t).
	\end{equation}
	
	撤销事件写入联盟链。随后母域将SMT中的设备状态修改为Revoked并生成新根：
	\begin{equation}
		Root_A^{new}=Update(Root_A,PID_j,Revoked).
	\end{equation}
	
	新根和IPFS CID被原子锚定到链上，同时将该设备加入硬件死名单。
	
	\subsection{硬件级不可洗白机制}
	软件身份和物理身份的分离是本文防洗白机制的核心。正常业务只暴露匿名PID，因此能够保护通信隐私；撤销时则将PUF派生的硬件关联值写入受保护的死名单。重新注册并不是简单地提交一个新的RealID，而是必须重新完成硬件证明。
	
	在攻击者只能修改软件层状态的威胁模型下，RealID、MAC和应用层密钥均不构成可靠的设备根身份。即使攻击者构造新的$PID'_j$，注册服务器仍能够从PUF重新获得同一硬件关联值，并与死名单进行比较。只有在硬件关联状态未被撤销的情况下，设备才可能重新获得注册凭证。
	
	设备重新注册时，必须重新完成PUF响应提取和注册证明。系统不直接把原始PUF响应上传链上，而是通过安全派生值构造硬件关联标识：
	\begin{equation}
		HID_j=H(K_j\parallel DomainSalt).
	\end{equation}
	
	若$HID_j$命中硬件死名单，则拒绝签发新的匿名凭证。由于攻击者只能改变软件层$RealID$、MAC或随机数而不能改变PUF根秘密，因此重新生成$PID'_j$不能绕过撤销状态。
	
	\section{密码与状态管理实例化}
	本文采用混合密码架构的原因在于不同密码原语具有不同的安全生命周期和计算特征。长期注册信任需要考虑后量子攻击，因此优先采用后量子签名；短时在线消息强调通信尺寸和验证效率，因此可采用BLS等短签名机制；状态证明需要低空间开销和可验证更新，因此采用SMT；高频集合筛选则使用Bloom Filter。
	
	这种组合并不意味着所有密码算法同时参与每一次认证。相反，每个算法仅在其最适合的阶段出现。这样可以避免协议描述中“算法堆叠”的问题，也使后续实验能够针对每个组件独立评估。
	
	\subsection{后量子注册凭证}
	需要区分“后量子签名”和“盲签名”两个概念。标准化ML-DSA能够提供后量子数字签名能力，但其接口本身并不自动提供盲签名属性。因此本文在系统层面将域级长期签名与匿名凭证签发接口分离。域级签名用于验证管理域和状态记录，匿名凭证则采用独立的格基盲签名构造。
	
	在实际实现中，可根据目标安全等级选择具体参数，并分别统计密钥、签名和证明尺寸。注册阶段允许较高成本，但仍需要考虑无人机存储能力，因此凭证材料应尽可能在注册服务器和边缘节点侧预处理。
	在线认证不重复执行注册阶段的高成本盲签名。注册凭证仅在设备加入系统或重新授权时更新。后量子签名组件可采用NIST标准化ML-DSA（原Dilithium）作为域级签名；匿名盲签名则采用独立的、具有相应安全证明的格基盲签名构造，避免概念混用。
	
	\subsection{BLS短签名}
	BLS的主要价值体现在签名尺寸较短以及支持聚合验证。在单请求低并发场景下，其优势并不一定显著，因此本文将其定位为在线短消息和高密度批量验证的可选实例化，而不是整个系统的长期根信任。多个无人机同时向同一边缘节点发送预认证消息时，可以将独立签名聚合后统一验证，从而减少验证接口调用次数。
	
	由于BLS依赖经典椭圆曲线配对，本文不把它作为后量子安全组件。长期信任仍由后量子签名和链上状态保护。这样的分层可以在保持在线性能的同时，避免把短期性能优化错误地表述为长期抗量子能力。
	BLS用于预认证阶段的短消息认证以及高密度场景中的批量验证：
	\begin{equation}
		\sigma_j=H_1(M_j)^{sk_j}.
	\end{equation}
	多个请求可进行聚合验证。考虑到BLS基于椭圆曲线配对、不是后量子安全原语，本文只将其用于短时在线通信，而不承担长期根信任。
	
	\subsection{SMT状态证明}
	SMT将设备状态组织为可验证的键值集合。叶节点包含匿名PID、状态、有效期和版本号，内部节点通过哈希递归计算得到状态根。联盟链不需要保存整个树，只需要保存最新根及其版本。目标域获得证明路径后即可在本地验证设备状态，而无需信任发送证明的中间节点。
	
	当设备被撤销时，仅需要更新受影响叶节点及其路径上的内部节点。与完整列表广播相比，SMT能够将状态证明规模控制在与树高相关的范围内。对于边缘节点而言，这种证明尤其适合处理Bloom Filter命中的少量候选设备。
	SMT叶节点采用：
	\begin{equation}
		Leaf_j=H(PID_j\parallel Status_j\parallel Exp_j\parallel Ver_j).
	\end{equation}
	链上仅保存根哈希，目标域通过Merkle路径或相应零知识封装证明设备状态。状态发生变化时只需更新相关路径并重新锚定根。
	
	\subsection{Bloom Filter参数}
	Bloom Filter的工程参数需要根据撤销集合规模和允许的假阳性率共同确定。位数组过小会导致假阳性率快速上升，使大量合法设备被迫进入SMT核验；位数组过大则增加边缘节点内存占用和同步流量。因此，系统应根据历史撤销增长率定期重新估计$n$并调整过滤器容量。
	
	本文不把假阳性率视为安全漏洞，而将其视为性能参数。只要命中后的处理路径始终进入精确状态验证，Bloom Filter的假阳性只会增加计算开销，不会改变最终授权结果。实验阶段应同时报告FPR、平均查询时间、内存占用和SMT二次验证比例。
	对于$n$个撤销设备、目标假阳性率$p$，Bloom Filter位数组长度可近似取：
	\begin{equation}
		m=-\frac{n\ln p}{(\ln2)^2},
	\end{equation}
	哈希函数数量为：
	\begin{equation}
		k=\frac{m}{n}\ln2.
	\end{equation}
	实际部署时根据撤销规模定期扩容或采用分代Bloom Filter。Bloom Filter只作为快速筛选结构，不替代链上状态证明。
	
	\subsection{IPFS与联盟链}
	区块链适合保存短小、需要一致性和不可篡改性的状态摘要，而不适合长期保存大量审计原文和完整状态树。因此本文采用“链上锚定、链下存储”的方式。IPFS保存完整SMT、历史状态和经过授权的审计证据，联盟链只保存Root、CID、Version和时间戳等摘要信息。
	
	目标域获取状态时首先验证链上根和版本，再根据CID获取对应链下对象。即使IPFS节点提供恶意或过期数据，目标域也能够通过根哈希发现不一致。这样可以在不显著增加链上负担的情况下保留完整审计能力。
	完整SMT、死名单分片和审计证据保存在IPFS。联盟链记录：
	\begin{equation}
		Record=\{Root,CID,Version,Timestamp,TxID\}.
	\end{equation}
	这样可以避免把大量叶节点和历史审计数据直接写入区块链，同时保持状态根和撤销事件的不可篡改性。
	
	\section{安全性分析}
	
	\subsection{双向认证}
	双向认证要求目标域不仅验证无人机的匿名凭证，也验证目标域自身的合法性。无人机接收目标域挑战和域级签名后，应验证域公钥是否属于联盟链登记的合法成员，并检查状态版本。只有通过双方验证后，双方才执行会话密钥建立。
	
	这种设计可以防止攻击者伪造一个“目标域”诱导无人机泄露认证材料。对于跨域飞行而言，目标域身份认证与无人机身份认证具有同等重要性，因为错误的目标域可能进一步导致控制权限或任务数据泄露。
	目标域只有在验证匿名凭证、挑战响应和最新SMT状态后才授予接入权限；无人机同时验证目标域的合法域证书和挑战信息，因此攻击者无法仅伪造PID获得接入权限。
	
	\subsection{抗重放}
	重放攻击在无人机无线链路中具有较高可行性。攻击者不需要破解签名，只需复制过去已经验证成功的认证消息并重新发送。本文通过时间戳、随机挑战和Nonce缓存形成三重约束。时间戳限制消息生命周期，目标域随机挑战保证不同正式认证会话之间的消息不可复用，Nonce缓存则阻止有效窗口内的重复发送。
	
	因此，即使历史消息的签名仍然有效，也不能在新的认证会话中直接使用。
	每次预认证和正式认证均包含随机Nonce与时间戳。Nonce缓存和时间窗口共同保证旧消息不能被重复接受。
	
	\subsection{匿名性与不可链接性}
	正常通信不直接传输RealID，而是使用匿名PID及必要的短期会话标识。目标域能够验证PID对应的凭证和状态，但普通网络观察者无法从单个报文恢复真实身份。对于需要更高不可链接性的应用，可在任务会话之间更新短期标识，而将长期硬件关联信息保留在受控注册侧。
	
	需要指出的是，匿名性并不等价于不可追责。本文采用受控映射机制，在发生严重安全事件并满足审计条件时，由授权管理实体完成真实身份追溯。这种设计兼顾正常业务隐私和事后安全治理。
	正常通信使用PID而非RealID。不同会话可以使用短期派生标识，避免长期轨迹直接关联。真实设备映射仅保留在授权管理域的受控审计环境中。
	
	\subsection{撤销正确性}
	撤销正确性需要同时考虑过滤器误报、状态同步延迟和旧凭证重用三个问题。Bloom Filter假阳性通过SMT二次验证消除；状态同步通过Root和Version进行检测；旧凭证则因为状态叶已经变为Revoked而无法通过正式认证。
	
	因此，系统允许存在“暂时不知道”的状态，但不允许在无法确认最新状态时默认授予高权限。对于版本落后的目标域，应暂停授权并请求状态同步，而不是继续使用旧状态。
	Bloom Filter命中不会直接导致拒绝，而是触发SMT精确核验。因此Bloom Filter假阳性不会导致合法设备被错误永久撤销。正式撤销由联盟链记录并由最新SMT根体现。
	
	\subsection{不可洗白性}
	不可洗白性建立在硬件根不可轻易复制这一假设之上。攻击者可以修改应用软件中的身份信息，却不能在不具备PUF物理特性的情况下生成与原设备一致的硬件派生秘密。注册服务器通过PUF关联值识别物理设备，因此新的软件PID不会自动产生新的信任根。
	
	需要注意的是，该安全性质依赖PUF实现质量。若实际硬件PUF存在可建模、可克隆或响应泄露问题，则防洗白能力会下降。因此，后续原型实验需要从稳定性、唯一性和攻击抵抗能力三个方面验证PUF实现。
	攻击者更换软件ID或重新生成随机数只能获得新的软件层标识，无法改变PUF根秘密。硬件死名单检查在重新注册之前执行，因此已撤销物理设备不能通过重新注册获得新的匿名凭证。
	
	\subsection{抗量子安全边界}
	本文采用“安全生命周期分层”的抗量子策略。长期存在的注册凭证、域级信任和链上身份签名优先采用后量子密码；短时会话和性能敏感的在线消息可以继续使用经过严格边界限定的经典高效机制。这样做的目标不是声称所有组件都具有相同的抗量子属性，而是避免经典密码成为长期身份根。
	
	在未来完全迁移到后量子在线密码后，BLS等组件可以替换为后量子聚合或短签名方案，而无需改变认证前移、过滤链、SMT和撤销的总体架构。
	长期域级信任和注册凭证采用后量子密码组件。BLS只用于短时在线消息认证，因此系统不将BLS的经典安全性作为长期根信任假设。
	
	\subsection{过滤链安全性}
	过滤链可以视为一个逐步增加计算成本的安全漏斗。第一层面对最低成本的明显异常，第二层处理已知状态异常，第三层处理行为风险，第四层处理撤销集合，最后才进入正式密码认证。设每一层的通过率分别为$p_1,p_2,p_3,p_4$，则进入正式认证的期望比例可近似表示为：
	\begin{equation}
		\rho=p_1p_2p_3p_4.
	\end{equation}
	该关系不意味着各层统计独立，而是用于说明过滤结构的资源缩减效果。实际实验应通过真实或仿真流量直接测量$\rho$，并分析不同恶意比例下各层的贡献。
	
	与此同时，过滤器不能改变合法设备最终获得授权的条件。所有“通过”结果仍需进入正式认证，所有“疑似撤销”结果仍需经过SMT确认，从而保证过滤层只是安全决策的前置优化，而不是降低认证标准。
	四层过滤形成递进式成本结构：
	\begin{equation}
		C_1<C_2<C_3<C_4<C_{auth}.
	\end{equation}
	大量低质量请求在前两层被丢弃；行为异常请求在第三层被分流；已知撤销设备在第四层快速命中；只有剩余请求进入正式密码认证，从而降低高成本认证接口暴露于请求洪泛攻击的程度。
	
	\section{性能评价方案}
	性能评价应覆盖通信、计算、存储和安全检测四个方面。仅报告一次认证时延无法体现分层过滤的实际价值，因此本文将比较不同过滤配置下的请求处理路径，并统计每层丢弃比例、CPU时间和内存占用。同时，将注册阶段和在线阶段分开评价，以验证“重型注册、轻量在线”的设计假设。
	
	实验应至少设置合法请求、外部伪造请求、重放请求、过期凭证请求和合法凭证恶意行为五类流量。这样可以分别验证四级过滤对不同攻击类型的适用性，而不是仅通过随机恶意流量获得一个综合准确率。
	
	由于完整系统尚处于协议实现阶段，本文不预先虚构实验结果。性能评价应从以下维度开展。
	
	\subsection{认证时延}
	认证前移的评价不能只比较密码运算时间，还应考虑边界触发时刻。对于每一次飞行任务，记录无人机进入触发区的时间$t_0$、预认证完成时间$t_1$以及实际越界时间$t_b$。若正式授权完成时间$t_2\le t_b$，则认为该任务满足无缝跨域条件。
	
	同时统计传统越界后认证的$T_{post}$与本文方案的$T_{pre}$。除了平均值，还应报告P95和P99时延，因为低空高并发系统更容易受到尾部延迟影响。
	比较传统“越界后认证”和本文“认证前移”机制，统计从进入边界缓冲区到正式授权完成的总时延：
	\begin{equation}
		T_{total}=T_{trigger}+T_{pre}+T_{filter}+T_{auth}.
	\end{equation}
	
	\subsection{过滤效果}
	过滤实验需要分别统计Precision、Recall、FPR和FNR。对于过滤系统而言，Precision反映被标记为恶意的请求中真正恶意请求的比例，Recall反映恶意请求被识别的程度；FPR用于评估合法设备被误送入审查路径的程度，FNR则关系到恶意请求漏过过滤链的风险。
	
	需要强调，过滤层不是最终授权模块，因此适度的假阳性是可以接受的，只要后续SMT和正式认证能够恢复最终判定。真正需要重点观察的是过滤后进入正式认证的请求数量以及恶意请求的最终漏检率。
	分别测试无过滤、单层过滤、两层过滤、三层过滤和完整四层过滤。主要指标为Precision、Recall、False Positive Rate、False Negative Rate和平均处理时间。论文不宣称过滤器具有100\%准确率，而通过实验测量实际误报和漏报。
	
	\subsection{Bloom Filter实验}
	Bloom Filter实验应在不同撤销规模和目标FPR下重复进行。随着撤销设备数量增加，如果保持固定的位数组长度，假阳性率将逐步上升；因此实验需要比较固定容量与动态扩容两种策略。
	
	除查询时间外，还应统计因假阳性进入SMT核验的请求比例。该指标能够直接反映Bloom Filter对后端状态证明模块的实际压力，比单独报告理论FPR更能说明方案的工程价值。
	改变撤销集合规模$n$、位数组$m$和哈希数量$k$，测量内存占用、查询时间和假阳性率，并验证假阳性只增加SMT核验次数而不会导致错误放行。
	
	\subsection{并发接入实验}
	并发实验重点考察边缘节点在请求洪泛和正常业务混合条件下的稳定性。逐步增加无人机数量和请求速率，记录过滤入口吞吐量、正式认证吞吐量、CPU占用、内存占用和平均响应时间。当攻击流量增加时，理想情况下前三层过滤应使高成本正式认证模块的负载增长速度低于入口请求增长速度。
	
	此外，应比较有无过滤链时区块链和SMT查询次数的变化，以验证过滤层是否真正降低了后端可信状态模块的压力。
	设置不同无人机数量，比较预认证请求吞吐量、正式认证请求数量和边缘节点CPU占用。重点验证过滤链是否能够减少进入高成本认证模块的请求比例。
	
	\subsection{注册与在线认证成本}
	注册成本和在线成本必须分别报告。注册阶段可以接受较长耗时，但应控制在设备入网流程可接受范围内；在线阶段则需要满足实时跨域切换要求。实验中应记录PUF提取、匿名凭证签发、凭证验证、BLS验证、SMT验证等操作的单次耗时。
	
	最终应给出单位时间内可支持的认证次数以及每次跨域任务的平均计算成本，而不是仅比较某一个密码算法的运行时间。
	分别统计注册阶段PUF提取、匿名凭证生成和服务器侧后量子操作成本，以及在线预认证、过滤和正式认证成本。预期系统以较高的低频注册成本换取较低的高频在线成本。
	
	\section{讨论}
	本文方案在设计时有意避免将所有安全问题交给单一密码算法解决，但这种分层结构也带来了参数选择和系统部署方面的新问题。首先，信任评分需要合理的行为特征和阈值，否则可能出现过度拒绝或检测灵敏度不足。其次，Bloom Filter需要随着撤销集合变化进行容量管理。再次，SMT状态和IPFS对象之间需要稳定的版本映射，以保证边缘节点不会长期使用过期状态。
	
	另外，PUF虽然能够提供物理设备绑定，但不同硬件平台的PUF稳定性差异较大。本文当前的协议设计以存在可靠PUF实现为前提，具体硬件实现需要通过温度、电压、老化和重复采样实验进行验证。对于没有硬件PUF的无人机，也可以使用安全芯片提供等价的不可导出设备根，但其安全属性需要单独分析。
	
	从部署角度看，本文更适合采用“边缘节点执行高频过滤、管理域执行授权决策、联盟链维护可信状态”的分布式架构。这样可以避免把所有任务集中到区块链节点，也可以降低无人机自身的在线计算压力。
	
	本文方案存在三个需要在原型实验中进一步验证的问题。第一，格基盲签名的具体参数需要依据最终安全等级确定，不能直接套用普通ML-DSA签名参数。第二，信任评分中的行为特征权重需要通过数据集或实测飞行日志标定，而不能人为选择后直接宣称最优。第三，PUF的稳定性和环境鲁棒性需要通过温度、电压和老化实验验证。
	
	此外，弹性边界中的$T_{auth}$和$T_{jit}$应由实际网络测试获得。本文只定义其协议位置和计算关系，不将未经测量的通信参数写成固定工程事实。
	
	\section{结论与展望}
	本文围绕低空无人机跨管理域飞行中的连续身份认证问题，构建了从设备注册到最终撤销的完整安全链路。与传统仅关注一次握手的认证协议不同，本文将飞行边界、请求过滤、状态证明、行为监测和设备撤销放入统一生命周期中考虑。注册阶段承担硬件绑定和后量子匿名凭证建立，认证前移机制负责提前准备跨域上下文，四级过滤链负责减少无效请求进入高成本认证模块，SMT和联盟链负责提供精确且可追溯的撤销状态，PUF死名单则进一步阻断被撤销物理设备重新获得身份。
	
	该设计的关键不在于简单增加安全组件，而在于重新安排安全操作的先后顺序：先以低成本规则和状态信息排除明显异常，再将有限计算资源用于需要密码学确认的请求；先在软件身份层保护匿名性，再在硬件根层维持撤销连续性。由此形成“轻量筛选、精确确认、可信撤销、硬件绑定”的安全闭环。
	
	后续研究将围绕三个方向展开。第一，完成格基匿名凭证的具体协议实例化，并给出不可伪造、不可链接和盲化安全性的形式化证明。第二，在NS-3或真实边缘平台上构建完整原型，测试不同无人机规模、网络延迟和恶意请求比例下的性能。第三，结合真实飞行日志训练和校准信任评分模型，并进一步研究在大规模撤销集合下的动态Bloom Filter和增量SMT更新策略。
	
	本文针对低空无人机跨管理域飞行中的认证滞后、恶意请求、动态撤销和身份洗白问题，提出一种认证前移与分层恶意请求过滤相结合的跨域认证方案。方案在低频注册阶段完成PUF硬件绑定和后量子匿名凭证建立，在高频在线阶段采用轻量预认证和四级过滤机制，并通过SMT提供精确状态证明。弹性边界使认证准备能够在实际越界前启动，协议/新鲜度、凭证状态、上下文信任和Bloom Filter四级过滤则形成由轻到重的请求处理链。设备被确认恶意后，联盟链记录不可逆撤销事件，SMT更新状态，PUF硬件死名单阻断重新注册，从而形成完整的“认证前移--请求过滤--正式认证--持续监测--精确撤销--硬件防洗白”闭环。
	
	后续工作将完成NS-3网络仿真和真实嵌入式平台测试，重点验证过滤链在不同恶意请求比例、无人机密度和无线信道条件下的实际性能；进一步研究格基匿名凭证的具体实例化、安全证明和通信尺寸；同时在ARM平台上测试PUF、后量子密码和在线认证流程的能耗，以确定方案在实际行业级无人机上的部署边界。
	
	\begin{thebibliography}{00}
		
		\bibitem{b1} L. Rui, X. Zhang, J. Yang, and L. Liu, ``Trust-Based Cross-Domain Batch Authentication Protocol for UAV Networks,'' in \emph{2025 11th International Conference on Computer and Communications (ICCC)}, 2025, pp. 1715--1720.
		
		\bibitem{b2} G. Qiao, P. Yang, T. Ye, and F. Han, ``A Blockchain-Based Cross-Domain Authentication and Flight Trajectory Privacy Protection Scheme for Unmanned Aerial Vehicle Networks,'' \emph{IEEE Internet of Things Journal}, vol. 13, no. 2, pp. 2155--2170, Jan. 2026.
		
		\bibitem{b3} R. Li, J. Cui, J. Zhang, L. Wei, H. Zhong, and D. He, ``Blockchain-Assisted Revocable Cross-Domain Authentication for Vehicular Ad-Hoc Networks,'' \emph{IEEE Transactions on Dependable and Secure Computing}, vol. 22, no. 5, pp. 4593--4607, 2025.
		
		\bibitem{b4} N. Kang, Z. Ning, S. Zhang, S. Rehman, and M. Waqas, ``Identity-Based Edge Computing Anonymous Authentication Protocol,'' \emph{Computers, Materials \& Continua}, vol. 74, no. 1, pp. 2005--2018, 2023.
		
		\bibitem{b5} Z. Cao, X. Wen, S. Ai, W. Shang, and S. Huan, ``A decentralized authentication scheme for smart factory based on blockchain,'' \emph{Scientific Reports}, vol. 14, no. 1, p. 24640, 2024.
		
		\bibitem{b6} K. Xue, X. Luo, Y. Ma, J. Li, J. Liu, and D. S. L. Wei, ``A Distributed Authentication Scheme Based on Smart Contract for Roaming Service in Mobile Vehicular Networks,'' \emph{IEEE Transactions on Vehicular Technology}, vol. 71, no. 5, pp. 5284--5297, 2022.
		
		\bibitem{b7} P. Gu, ``An Efficient Blockchain-based Cross-domain Authentication and Secure Certificate Revocation Scheme,'' in \emph{2020 IEEE 6th International Conference on Computer and Communications (ICCC)}, 2020, pp. 1651--1655.
		
		\bibitem{b8} Y. Sun, R. Lu, X. Lin, X. S. Shen, and J. Su, ``An Efficient Pseudonymous Authentication Scheme With Strong Privacy Preservation for Vehicular Communications,'' \emph{IEEE Transactions on Vehicular Technology}, vol. 59, no. 7, pp. 3589--3603, 2010.
		
		\bibitem{b9} C. Lin, D. He, X. Huang, N. Kumar, and K. R. Choo, ``BCPPA: A Blockchain-Based Conditional Privacy-Preserving Authentication Protocol for Vehicular Ad Hoc Networks,'' \emph{IEEE Transactions on Intelligent Transportation Systems}, vol. 22, no. 12, pp. 7408--7420, 2021.
		
		\bibitem{b10} M. Shen, H. Liu, L. Zhu, K. Xu, H. Yu, X. Du, and M. Guizani, ``Blockchain-Assisted Secure Device Authentication for Cross-Domain Industrial IoT,'' \emph{IEEE Journal on Selected Areas in Communications}, vol. 38, no. 5, pp. 942--954, 2020.
		
		\bibitem{b11} Z. Li, X. Liu, N. Zhang, and Z. Wei, ``Blockchain-Based Certificateless Cross-Domain Authentication Scheme in the Industrial Internet of Things,'' \emph{Computers, Materials \& Continua}, vol. 80, no. 3, pp. 3835--3856, 2024.
		
		\bibitem{b12} J. Dong, G. Xu, M. Chuang, J. Liu, and U. G. O. Cliff, ``Blockchain-Based Certificate-Free Cross-Domain Authentication Mechanism for Industrial Internet,'' \emph{IEEE Internet of Things Journal}, vol. 11, no. 2, pp. 3316--3329, 2024.
		
		\bibitem{b13} F. Wang, J. Cui, Q. Zhang, D. He, C. Gu, and H. Zhong, ``Blockchain-Based Lightweight Message Authentication for Edge-Assisted Cross-Domain Industrial Internet of Things,'' \emph{IEEE Transactions on Dependable and Secure Computing}, vol. 21, no. 4, pp. 1587--1601, 2024.
		
		\bibitem{b15} H. Zhong, C. Gu, Q. Zhang, J. Cui, C. Gu, and D. He, ``Conditional privacy-preserving message authentication scheme for cross-domain Industrial Internet of Things,'' \emph{Ad Hoc Networks}, vol. 144, p. 103137, 2023.
		
		\bibitem{b16} J. Cui, N. Liu, Q. Zhang, D. He, C. Gu, and H. Zhong, ``Efficient and Anonymous Cross-Domain Authentication for IIoT Based on Blockchain,'' \emph{IEEE Transactions on Network Science and Engineering}, vol. 10, no. 3, pp. 899--910, 2023.
		
		\bibitem{b17} M. Zeng, J. Cui, Q. Zhang, H. Zhong, and D. He, ``Efficient Revocable Cross-Domain Anonymous Authentication Scheme for IIoT,'' \emph{IEEE Transactions on Information Forensics and Security}, vol. 20, pp. 996--1010, 2025.
		
		\bibitem{b18} Y. Chen, G. Dong, J. Bai, Y. Hao, F. Li, and H. Peng, ``A Trust Enhancement Scheme for Cross Domain Authentication of PKI system,'' in \emph{2019 International Conference on Cyber-Enabled Distributed Computing and Knowledge Discovery}, 2019, pp. 364--371.
		
		\bibitem{b19} Y. Wu, Z. Jia, J. Zhao, H. Wang, Q. Wu, and Z. Han, ``A Blockchain-Based Trust Framework for Resilient Cross-Domain UAV Service Orchestration,'' arXiv preprint arXiv:2603.00456, 2026.
		
		\bibitem{b20} S. Liang, B. Duo, J. Ning, and R. Zhou, ``A Lightweight Authentication and Key Agreement Scheme for Cross-Domain UAV Networks Based on Extended Chebyshev Chaotic Map and Hash Chain,'' in \emph{2025 4th International Conference on Cloud Computing, Big Data Application and Software Engineering}, 2025.
		
		\bibitem{b21} L. Feng, C. Xu, M. Liu, F. Huang, and Y. Yu, ``TEBP-UAVs: A trusted and efficient blockchain-based protocol for cross-domain authentication in UAVs,'' \emph{Computer Networks}, vol. 282, p. 112243, 2026.
		
	\end{thebibliography}
	
\end{document}
